% Part: computability
% Chapter: computability-theory
% Section: coding-computations

\documentclass[../../../include/open-logic-section]{subfiles}

\begin{document}

\olfileid{cmp}{thy}{cod}
\olsection{গণনার সংকেতায়ন}

প্রত্যেক গণনা-মডেলে নিচের কাজগুলি করা সম্ভব:
\begin{enumerate}
\item গণনসাধ্য অপেক্ষকের \emph{সংজ্ঞাগুলি} নিয়মমাফিক বর্ণনা করা। যেমন,
  টুরিং যন্ত্রের নির্দেশন, পুনরাবৃত্ত সংজ্ঞা বা কোনো প্রোগ্রামিং ভাষার
  প্রোগ্রামকে এই সংজ্ঞাগুলি দেওয়ার উপায় হিসেবে ভাবা যায়।
\item কোনো নির্দিষ্ট নিবেশে কোনো সংজ্ঞা-নির্ধারিত অপেক্ষকের গণনার পূর্ণ
  নথি বর্ণনা করা। যেমন, গণনার প্রত্যেক ধাপের কনফিগারেশনের অনুক্রম
  (যন্ত্রটির দশা ও ফিতার বিষয়বস্তু) দিয়ে টুরিং যন্ত্রের গণনা বর্ণনা করা
  যায়।
\item গণনার সম্ভাব্য কোনো নথি সত্যিই কোনো নির্দিষ্ট নিবেশে নির্দিষ্ট
  সংজ্ঞার গণনসাধ্য অপেক্ষকটির গণনার নথি কি না পরীক্ষা করা (যে নিবেশে
  অপেক্ষকটি সংজ্ঞায়িত, অর্থাৎ গণনাটি থামে)।
\item গণনার পূর্ণ নথির এমন বর্ণনা থেকে নির্দিষ্ট নিবেশে অপেক্ষকটির মান
  বার করা। যেমন, থেমে যাওয়া টুরিং-যন্ত্রের গণনার একেবারে শেষ ধাপে ফিতার
  বিষয়বস্তুই মানটি।
\end{enumerate}

সংকেতায়নের সাহায্যে গণনসাধ্য অপেক্ষকের প্রত্যেক বর্ণনাকে এমনভাবে একটি
সংখ্যাসূচক \emph{সূচক} দেওয়া যায়, যাতে সূচক থেকে গণনসাধ্য উপায়ে
নির্দেশগুলি পুনরুদ্ধার করা যায়। একইভাবে, কোনো গণনার পূর্ণ নথিকেও একটি
মাত্র সংখ্যা দিয়ে সংকেতায়িত করা যায়। তখন প্রাপ্ত পাটিগাণিতিক সম্বন্ধ
“নিবেশ~$x$-এর জন্য সূচক~$e$-এর অপেক্ষকটির গণনার নথির সংকেত হলো~$s$”
এবং অপেক্ষক “সংকেত~$s$-এর গণনা-অনুক্রমের নির্গম”—দুটিই গণনসাধ্য; আসলে
দুটিই আদিম পুনরাবৃত্ত।

এই মৌলিক তথ্যটি অত্যন্ত শক্তিশালী। নির্বাচিত গণনা-মডেলের উপর নির্ভর না
করেই এর সাহায্যে গণনসাধ্যতা সম্বন্ধে কয়েকটি চমকপ্রদ ও গুরুত্বপূর্ণ ফল
প্রমাণ করা যায়।

\end{document}
